\section{Estado del arte}

La implementación de controles predictivos avanzados como el control predictivo basado en modelos (MPC) se ve 
limitada por el alto costo computacional asociado a los modelos rigurosos que describen procesos químicos complejos. 
Al requerir cálculos numéricos intensivos, la subrogación de modelos surge como una técnica eficiente para representar 
el comportamiento del sistema mediante aproximaciones matemáticas. A través de metamodelos, es posible establecer 
relaciones entre variables de entrada y salida con el objetivo de mantener la integridad del control mientras se reduce 
la carga computacional \cite{ref6}. Diversos estudios han explorado la subrogación de modelos con el propósito de 
integrarlos en sistemas de control predictivo.

\subsection{Modelos ARX y ARMAX}
Los modelos ARX (AutoRegressive with eXogenous inputs) constituyen una clase de modelos lineales de identificación de 
sistemas en tiempo discreto que describen la salida actual del sistema como una combinación lineal de sus salidas pasadas
y de las entradas actuales y pasadas \cite{ref11}. Su estructura general se expresa como:

\begin{equation}
A(q) \bar{y}(k) = B(q) \bar{u}(k-n_k) + \bar{e}(k)
\label{eq:arx_general}
\end{equation}

\begin{equation}
A(q) = a_1 + a_2q^{-1} + \cdots + a_{n} q^{-n_a}
\label{eq:arx_A}
\end{equation}

\begin{equation}
B(q) = b_1 + b_2q^{-1} + \cdots + b_{n} q^{-n_b+1}
\label{eq:arx_B}
\end{equation}

En las ecuaciones \ref{eq:arx_general} \ref{eq:arx_B}, $y(k)$ representa la salida del sistema en el instante $k$, $u(k)$ 
la entrada manipulada, $A(q)$ y $B(q)$ los polinomios que contienen los parámetros del modelo, $n_a$ y $n_b$ el número de 
polos y ceros del sistema, respectivamente, y $n_k$ el retraso del modelo. El término $e(k)$ corresponde al ruido blanco 
que captura las perturbaciones no modeladas y $q^{-1}$ es el operador de retraso. Los parámetros del modelo se estiman típicamente mediante el método de mínimos cuadrados ordinarios, lo que garantiza una 
solución analítica cerrada \cite{ref11}. Una ventaja importante de los modelos ARX en el contexto del control MPC es su 
compatibilidad directa con herramientas computacionales de identificación de sistemas, lo que permite su integración 
inmediata en esquemas de control sin necesidad de transformaciones adicionales. Estudios recientes han aplicado modelos ARX en la identificación y control de columnas de destilación, reportando  desempeños adecuados incluso cuando el proceso presenta comportamientos ligeramente no lineales \cite{ref12}. Sin embargo,  su desempeño puede verse limitado en presencia de fuertes no linealidades, situación en la que estructuras como los modelos NARX o modelos de mayor complejidad pueden ofrecer mejores resultados \cite{ref13}. No obstante, cuando el sistema presenta un nivel de ruido significativo, resulta conveniente utilizar modelos ARMAX, los
cuales incorporan un término adicional que modela explícitamente la dinámica del error:

\begin{equation}
C(q) = c_1 + c_2q^{-1} + \cdots + c_{n_c} q^{-n_c}
\label{eq:armax_C}
\end{equation}

\subsection{Representación en espacio de estados}

La representación en espacio de estados describe la dinámica de un sistema mediante un conjunto de ecuaciones diferenciales de primer orden que capturan su evolución temporal a través de un conjunto de variables internas. Estas variables se agrupan en el vector de estados, el cual contiene la información mínima necesaria para caracterizar el comportamiento dinámico del proceso y establecer la relación entre las entradas y las salidas del sistema. De manera general, esta representación se expresa mediante las ecuaciones \ref{eq:statespace_state} y \ref{eq:statespace_output}

\begin{equation}
\dot{x}(t) = A x(t) + B u(t)
\label{eq:statespace_state}
\end{equation}

\begin{equation}
y(t) = C x(t)
\label{eq:statespace_output}
\end{equation}

Donde $x(t)$ representa un vector con las derivadas de primer orden de las variables de estado, $x(t)$ es el vector de variables de estado, $u(t)$ es el vector con las variables de entrada, $y(t)$ es el vector con las variables de salida y $A$, $B$, y $C$ son las matrices de estado, entrada y salida respectivamente. Diversos estudios han aplicado esta formulación para modelar y controlar sistemas dinámicos complejos. Por ejemplo, utiliznado una representación en espacio de estados para convertir una simulación dinámica rigurosa en un modelo matemático simplificado se puede obtener la matriz de transferencia del sistema\cite{ref20}, lo que permite diseñar un controlador óptimo mediante la técnica de realimentación de ganancia, facilitando el diseño del esquema de control y mejorando el desempeño dinámico del proceso. En otra investigación se empleo el espacio de estados en la identificación de sistemas dinámicos no lineales basada en datos mediante autocodificadores profundos, un tipo de redes neuronales\cite{ref21}. En este enfoque, el autocodificador aprende una representación latente del sistema, es decir, un conjunto de variables internas no observables directamente, pero capaces de capturar la dinámica del proceso, que actúan como vectores de estado. Los autores validaron esta aproximación mediante simulaciones en lazo abierto y posteriormente evaluaron su desempeño dentro de un esquema de control predictivo basado en modelos (MPC) en lazo cerrado. Otros estudios han desarrollado un esquema de MPC para el control de procesos por lotes afectados por perturbaciones desconocidas y fallas en actuadores. En este trabajo emplearon una representación no mínima del espacio de estados, es decir, una formulación que incorpora más estados de los estrictamente necesarios para describir la dinámica de entrada–salida del proceso \cite{ref22}. Esta ampliación permitió mejorar el desempeño en el lazo cerrado, ya que el sistema logró compensar simultáneamente las variaciones en las entradas y los errores de seguimiento.

\subsection{Modelo Hammerstein--Wiener}

Los modelos Hammerstein–Wiener constituyen una estructura ampliamente utilizada para la identificación de sistemas no lineales. En este enfoque, la dinámica del proceso se describe mediante un bloque lineal formulado en espacio de estados, mientras que las variables de entrada y/o salida se transforman mediante funciones no lineales estáticas, es decir, funciones sin dinámica ni memoria. En particular, el bloque Hammerstein introduce una no linealidad estática en la entrada antes del sistema lineal, mientras que el bloque Wiener aplica una transformación no lineal a la salida del sistema lineal. Esta arquitectura permite representar comportamientos no lineales al combinar un modelo dinámico lineal, responsable de describir la evolución temporal del sistema, con transformaciones estáticas que pueden capturar fenómenos como saturaciones, zonas muertas o ganancias dependientes del estado \cite{ref23}. Debido a esta estructura híbrida, los modelos Hammerstein–Wiener resultan especialmente adecuados para procesos químicos, ya que mejoran la capacidad de representar relaciones no lineales entre las variables de entrada y salida del sistema.

\subsection{Resumen de las estrategias de subrogación}

La elección de la técnica de subrogación de modelos para el control y la optimización de procesos químicos no lineales debe basarse en las características específicas del sistema que se desea modelar. Los modelos ARX y ARMAX son adecuados cuando el proceso presenta una dinámica lineal o ligeramente no lineal, lo que permite una representación eficiente con un costo computacional bajo. Sin embargo, su aplicabilidad se ve limitada cuando el sistema presenta no linealidades complejas o perturbaciones significativas. En estos casos, las máquinas de soporte vectorial (SVM) ofrecen una mayor robustez, ya que son capaces de manejar relaciones no lineales más complejas y de generalizar bien a partir de conjuntos de datos limitados. Las SVM se convierten en una opción poderosa para sistemas con alta incertidumbre y variabilidad en sus dinámicas. A pesar de sus ventajas, su integración en esquemas de control predictivo, como el MPC, presenta desafíos debido a que los modelos no son directamente lineales y requieren técnicas adicionales para su implementación en tiempo real. Por último, la representación en espacio de estados es una opción más robusta para procesos con dinámicas altamente no lineales o para sistemas que involucran interacciones complejas entre múltiples variables. Esta técnica permite capturar de manera precisa las interacciones internas y externas del sistema, pero a costa de un mayor costo computacional, debido a la necesidad de resolver sistemas de ecuaciones diferenciales. Por todo lo anterior, se busca comparar los diferentes métodos para determinar sus diferencias en el caso propuesto. 
En la Tabla \ref{tab:comparacion_modelos} se presenta una comparación general entre las principales estrategias de modelado utilizadas en 
procesos químicos.

\begin{table}[ht]
\centering
\small
\caption{Comparación de estrategias de subrogación para control de procesos}
\label{tab:comparacion_modelos}
\renewcommand{\arraystretch}{1.3}
\setlength{\tabcolsep}{6pt}
\begin{tabularx}{\textwidth}{l l l l l X}
\hline
Modelo & Tipo & No linealidad & Complejidad & Integración MPC & Aplicaciones \\
\hline
Espacio de estados & Lineal dinámico & Baja--media & Media & Directa & Control multivariable \\
\hline
ARX / ARMAX & Lineal dinámico & Media & Media & Directa & Identificación de sistemas \\
\hline
SVM & No lineal & Alta & Alta & Requiere MPC no lineal & Sistemas no lineales \\
\hline
Hammerstein--Wiener & Híbrido & Media--alta & Media & Compatible & Procesos químicos \\
\hline
\end{tabularx}
\end{table}



\subsection{Métodos de modelamiento de columnas de destilación}

El modelamiento de columnas de destilación se basa comúnmente en el modelo MESH (Mass, Equilibrium, Summation and Heat), el cual asume 
equilibrio termodinámico entre las fases vapor y líquido en cada etapa de la columna \cite{ref10}. El número total de ecuaciones algebraicas que describen una columna de destilación en un instante dado puede calcularse mediante:

\begin{equation}
N_{ec} = N (2C + 3)
\label{eq:num_ec}
\end{equation}

donde $N$ representa el número de etapas teóricas de la columna y $C$ el número de componentes químicos presentes en el sistema. En simulaciones dinámicas, este sistema de ecuaciones se resuelve para cada intervalo de tiempo considerado en la discretización del 
proceso. Por lo tanto, el número total de ecuaciones que debe resolver el simulador se expresa como:

\begin{equation}
N_{et} = N_{ec} \times N_{i}
\label{eq:num_et}
\end{equation}

donde $N_i$ representa el número de intervalos de tiempo utilizados en la simulación dinámica. El esquema general del proceso de destilación considerado en este trabajo se presenta en el Anexo correspondiente.